\textbf{Сверка хеша кода (скрипт на стр. 1).}
Шесть символов над листингом помогают искать опечатки.
Сохранить скрипт как \texttt{hash.sh}, \emph{только код листинга} —
как \texttt{snippet.cpp}: без main, общих include и pragma once.
Зависимости над листингом не входят в хеш. В терминале:
\begin{lstlisting}[language={}]
sh hash.sh < snippet.cpp
# Или нужные строки из решения:
sed -n '20,70p' solution.cpp | sh hash.sh
\end{lstlisting}
20..70 заменить своими границами. Вывод сравнить с хешем в PDF:
для \texttt{RomanNumerals.h} это \texttt{d47dc3}.
Сверять \emph{до} изменения имён, констант и выбора вариантов.
Комментарии и пробелы игнорируются. Ошибки cpp сначала исправить.
Совпадение не заменяет тесты: хеш короткий, пробелы удаляются даже
внутри строковых литералов. Подробнее: \texttt{doc/code-hash.md}.
